% ============================================================
% CHAPTER 7: Sprint 4 — CI/CD Integration, Docker Deployment and Demo
% (Following Tarek's sprint chapter pattern)
% ============================================================

\chapter{Sprint 4: CI/CD Integration, Docker Deployment and Demo}

\section{Introduction}

This chapter covers the final sprint (March 16 -- April 10, 2026), the longest at four weeks. During this iteration, we containerized the service with Docker, implemented the Jenkins CI/CD pipeline integration, set up the Gitea local Git server, and built the complete multi-developer offline demonstration. This sprint transformed the project from a standalone service into a fully integrated DevSecOps pipeline, proving that automated infrastructure security scanning can operate entirely offline without cloud dependencies. The chapter culminates with the Alice and Bob demonstration scenario that showcases the end-to-end workflow.

% ============================================================
\section{Sprint 4 Objectives}
% ============================================================

The objectives for Sprint 4 were:

\begin{enumerate}
    \item Containerize the service with Docker and Docker Compose.
    \item Create the Jenkinsfile for automated CI/CD pipeline integration.
    \item Set up Jenkins on localhost:8080 and Gitea on localhost:3000.
    \item Implement the multi-developer demonstration (Alice vs.\ Bob scenario).
    \item Create CI/CD integration examples for Jenkins, GitLab CI, and GitHub Actions.
    \item Validate the complete workflow from code push to pipeline decision.
\end{enumerate}

% ============================================================
\section{Sprint 4 Backlog}
% ============================================================

\begin{table}[H]
\centering
\caption{Sprint 4 — Product Backlog}
\label{tab:sprint4-backlog}
\begin{tabular}{|p{0.8cm}|p{8.5cm}|c|c|}
\hline
\textbf{ID} & \textbf{User Story} & \textbf{Priority} & \textbf{Status} \\
\hline
US-17 & As a DevOps engineer, I can deploy the service using Docker and docker-compose with a single command. & Medium & Done \\
\hline
US-18 & As a DevOps engineer, I can integrate the service into a Jenkins pipeline to automatically scan code on every push. & High & Done \\
\hline
US-19 & As a DevOps engineer, I can configure the pipeline to FAIL when security violations are detected. & High & Done \\
\hline
US-20 & As a developer (Alice), I can see my pipeline FAIL when I push insecure Terraform code. & High & Done \\
\hline
US-21 & As a developer (Bob), I can see my pipeline PASS when I push secure Terraform code. & High & Done \\
\hline
US-22 & As an administrator, I can run the complete demonstration environment offline. & Medium & Done \\
\hline
\end{tabular}
\end{table}

% ============================================================
\section{Use Case Refinement}
% ============================================================

\subsection{Refined Use Case: End-to-End CI/CD Pipeline}

% TODO: INSERT USE CASE DIAGRAM FOR SPRINT 4
% Show the full CI/CD flow:
% Developer → Gitea → Jenkins → Policy Service → ALLOW/BLOCK → Deploy or Fail
% \begin{figure}[H]
%     \centering
%     \includegraphics[width=0.95\textwidth]{figures/sprint4-usecase.png}
%     \caption{Sprint 4 — Refined use case: end-to-end CI/CD pipeline}
%     \label{fig:sprint4-usecase}
% \end{figure}

\textit{[Figure: Sprint 4 Use Case Diagram — create in draw.io]}

\vspace{0.3cm}

\begin{table}[H]
\centering
\caption{Textual description — Automated pipeline security scan}
\label{tab:sprint4-usecase-text}
\begin{tabular}{|p{3.5cm}|p{9.5cm}|}
\hline
\textbf{Use Case} & Automated Pipeline Security Scan on Git Push \\
\hline
\textbf{Primary Actor} & Developer (Alice or Bob) \\
\hline
\textbf{Secondary Actors} & Gitea (Git server), Jenkins (CI/CD), Policy Service \\
\hline
\textbf{Precondition} & All three services are running: Policy Service (port 8000), Jenkins (port 8080), Gitea (port 3000). A repository is configured with a Jenkinsfile. \\
\hline
\textbf{Postcondition} & If insecure: pipeline FAILS and deployment is blocked. If secure: pipeline PASSES and deployment proceeds. Dashboard is updated with the scan results. \\
\hline
\textbf{Nominal Scenario} &
\begin{enumerate}[nosep]
    \item The developer pushes infrastructure code to the Gitea repository.
    \item Gitea triggers a webhook to notify Jenkins of the new commit.
    \item Jenkins starts the pipeline defined in the Jenkinsfile.
    \item The pipeline reads the infrastructure files and sends each to the policy service via \texttt{POST /analyze}.
    \item The policy service evaluates all 100 policies and returns ALLOW or BLOCK.
    \item If BLOCK: Jenkins marks the build as FAILED and logs the violations.
    \item If ALLOW: Jenkins proceeds to the build and deploy stages.
    \item The dashboard updates with the new scan data.
\end{enumerate} \\
\hline
\end{tabular}
\end{table}

% ============================================================
\section{Sequence Diagram}
% ============================================================

% TODO: INSERT SEQUENCE DIAGRAM FOR SPRINT 4
% Full CI/CD flow with all 5 actors:
% Developer → Gitea → Jenkins → PolicyService → SQLite
% \begin{figure}[H]
%     \centering
%     \includegraphics[width=\textwidth]{figures/sprint4-sequence.png}
%     \caption{Sequence diagram — complete CI/CD pipeline flow}
%     \label{fig:sprint4-sequence}
% \end{figure}

\textit{[Figure: Sprint 4 Sequence Diagram — create in draw.io]}

% ============================================================
\section{Realization}
% ============================================================

\subsection{Docker Containerization}

The service is packaged as a Docker container using a production-ready Dockerfile that follows Docker security best practices:

\begin{lstlisting}[language=bash, caption={Dockerfile — production-ready container}, label={lst:dockerfile}]
FROM python:3.10-slim

LABEL maintainer="ridha"
LABEL description="DevSecOps Policy-as-Code Microservice"

WORKDIR /app
COPY requirements.txt .
RUN pip install --no-cache-dir --upgrade pip && \
    pip install --no-cache-dir -r requirements.txt
COPY . .

# Security: run as non-root user
RUN useradd -m -u 1000 appuser && \
    chown -R appuser:appuser /app
USER appuser

EXPOSE 8000

HEALTHCHECK --interval=30s --timeout=3s \
  --start-period=5s --retries=3 \
  CMD python -c "import requests; \
    requests.get('http://localhost:8000/health')" || exit 1

CMD ["python", "main.py"]
\end{lstlisting}

\noindent Notably, this Dockerfile follows the very best practices that our own policies enforce:
\begin{itemize}
    \item Uses a specific version tag (\texttt{python:3.10-slim}), not \texttt{latest}.
    \item Creates and runs as a non-root user (\texttt{appuser}).
    \item Includes a HEALTHCHECK instruction for container orchestration.
    \item Uses \texttt{--no-cache-dir} to minimize image size.
\end{itemize}

Docker Compose orchestrates the deployment:

\begin{lstlisting}[language=bash, caption={docker-compose.yml — single-command deployment}, label={lst:docker-compose}]
version: '3.8'
services:
  policy-service:
    build: .
    container_name: devsecops-policy-service
    ports:
      - "8000:8000"
    environment:
      - LOG_LEVEL=INFO
      - PYTHONUNBUFFERED=1
    volumes:
      - ./config:/app/config:ro
    restart: unless-stopped
    healthcheck:
      test: ["CMD", "curl", "-f", 
             "http://localhost:8000/health"]
      interval: 30s
      timeout: 10s
      retries: 3
\end{lstlisting}

\subsection{Jenkins Pipeline Integration}

The Jenkinsfile defines a multi-stage pipeline that automatically scans all infrastructure files in a repository. The key stages are:

\begin{enumerate}
    \item \textbf{Start Policy Service}: Launches the Docker container and waits for the health check.
    \item \textbf{Security Scan}: Iterates through all \texttt{.tf}, \texttt{Dockerfile}, and \texttt{.yaml} files, sending each to the \texttt{POST /analyze} endpoint.
    \item \textbf{Get Metrics}: Retrieves aggregated scan statistics for the build log.
    \item \textbf{Deploy}: Proceeds only if all security scans pass (no BLOCK decisions).
\end{enumerate}

\begin{lstlisting}[language=Groovy, caption={Jenkinsfile — security scan stage (simplified)}, label={lst:jenkinsfile}]
stage('Security Scan') {
    steps {
        script {
            def violations = []
            
            // Scan all Terraform files
            def terraformFiles = findFiles(glob: '**/*.tf')
            terraformFiles.each { file ->
                def content = readFile(file.path)
                def payload = JsonOutput.toJson([
                    code_type: 'terraform',
                    content: content
                ])
                
                def response = sh(
                    script: """curl -s -X POST \
                      http://localhost:8000/analyze \
                      -H 'Content-Type: application/json' \
                      -d '${payload}'""",
                    returnStdout: true
                ).trim()
                
                def result = readJSON text: response
                
                if (result.decision == 'BLOCK') {
                    violations.add([file: file.path, 
                        violations: result.violations])
                }
            }
            
            if (violations.size() > 0) {
                error("Security violations detected")
            }
        }
    }
}
\end{lstlisting}

\subsection{Gitea Setup}

Gitea runs locally on port 3000, providing a lightweight Git hosting service. The setup involves:

\begin{enumerate}
    \item Installing Gitea as a single binary on the development machine.
    \item Creating two user accounts: \textbf{alice} and \textbf{bob}.
    \item Creating a shared repository with the Jenkinsfile and infrastructure code.
    \item Configuring a webhook to notify Jenkins on push events.
\end{enumerate}

% TODO: INSERT SCREENSHOT OF GITEA REPOSITORY
% \begin{figure}[H]
%     \centering
%     \includegraphics[width=0.9\textwidth]{figures/sprint4-gitea.png}
%     \caption{Sprint 4 — Gitea repository with infrastructure code}
%     \label{fig:sprint4-gitea}
% \end{figure}

\textit{[Screenshot: Gitea repository at localhost:3000]}

\subsection{Alice vs.\ Bob Demonstration Scenario}

The demonstration showcases the complete DevSecOps workflow through two contrasting developer personas:

\vspace{0.3cm}
\noindent\textbf{Alice (Insecure Code):}

Alice pushes a Terraform file containing multiple security violations:
\begin{itemize}
    \item A public S3 bucket (\texttt{acl = "public-read"})
    \item An unencrypted RDS instance (\texttt{storage\_encrypted = false})
    \item A publicly accessible database (\texttt{publicly\_accessible = true})
\end{itemize}

\noindent\textbf{Result}: The policy service detects 19 violations $\rightarrow$ \textcolor{red}{\textbf{BLOCK}} $\rightarrow$ Jenkins pipeline \textbf{FAILS} $\rightarrow$ deployment is prevented.

\vspace{0.3cm}
\noindent\textbf{Bob (Secure Code):}

Bob pushes a properly configured Terraform file:
\begin{itemize}
    \item Private S3 bucket (\texttt{acl = "private"})
    \item Server-side encryption enabled (\texttt{sse\_algorithm = "AES256"})
    \item Compliance tags present (\texttt{Environment}, \texttt{Owner}, \texttt{CostCenter})
\end{itemize}

\noindent\textbf{Result}: The policy service finds 0 violations $\rightarrow$ \textcolor{green!50!black}{\textbf{ALLOW}} $\rightarrow$ Jenkins pipeline \textbf{PASSES} $\rightarrow$ deployment proceeds.

\vspace{0.3cm}
\noindent\textbf{Dashboard after both scans}:
\begin{itemize}
    \item Total scans: 2
    \item Blocked: 1 (Alice)
    \item Allowed: 1 (Bob)
    \item Block rate: 50\%
\end{itemize}

% TODO: INSERT ALICE vs BOB COMPARISON DIAGRAM
% \begin{figure}[H]
%     \centering
%     \includegraphics[width=0.9\textwidth]{figures/sprint4-alice-bob.png}
%     \caption{Sprint 4 — Alice (BLOCK) vs. Bob (ALLOW) comparison}
%     \label{fig:sprint4-alice-bob}
% \end{figure}

\textit{[Figure: Alice vs. Bob Comparison Diagram — create in draw.io]}

\subsection{CI/CD Integration Templates}

In addition to the Jenkins integration, we provided ready-to-use pipeline templates for two other popular CI/CD platforms:

\vspace{0.3cm}
\noindent\textbf{GitLab CI} (\texttt{.gitlab-ci.yml}): A job that scans Terraform files and fails the pipeline on violations.

\noindent\textbf{GitHub Actions} (\texttt{github-actions.yml}): A workflow triggered on pull requests that runs the policy service as a Docker container and comments scan results on the PR.

These templates demonstrate that the service integrates with any CI/CD platform that supports HTTP API calls.

% ============================================================
\section{Technologies Used}
% ============================================================

\vspace{0.5cm}
\begin{tabular}{| m{2.5cm} | m{10cm} |}
\hline

\textbf{Docker}~\cite{docker} &
OS-level containerization platform. The policy service runs as a Docker container with health checks, non-root user, and minimal image size.
\\
\hline

\textbf{Docker Compose}~\cite{dockercompose} &
Multi-container orchestration tool. Enables single-command deployment of the policy service with its configuration volume.
\\
\hline

\textbf{Jenkins 2.x}~\cite{jenkins} &
Open-source CI/CD automation server. Runs locally on port 8080 and executes the Jenkinsfile pipeline on code push events.
\\
\hline

\textbf{Gitea} &
Lightweight, self-hosted Git service written in Go. Runs locally on port 3000, providing repository management and webhook triggers for the offline demonstration.
\\
\hline

\textbf{Groovy} &
Scripting language used in Jenkinsfiles for pipeline logic, JSON parsing, and build decision control.
\\
\hline

\end{tabular}

% ============================================================
\section{Sprint 4 Interfaces}
% ============================================================

\subsection{Jenkins Pipeline — Failed Build (Alice)}

% TODO: INSERT SCREENSHOT OF JENKINS PIPELINE FAILURE
% \begin{figure}[H]
%     \centering
%     \includegraphics[width=0.9\textwidth]{figures/sprint4-jenkins-fail.png}
%     \caption{Sprint 4 — Jenkins pipeline FAIL — Alice's insecure code blocked}
%     \label{fig:sprint4-jenkins-fail}
% \end{figure}

\textit{[Screenshot: Jenkins pipeline with red FAIL stage]}

\subsection{Jenkins Pipeline — Successful Build (Bob)}

% TODO: INSERT SCREENSHOT OF JENKINS PIPELINE SUCCESS
% \begin{figure}[H]
%     \centering
%     \includegraphics[width=0.9\textwidth]{figures/sprint4-jenkins-pass.png}
%     \caption{Sprint 4 — Jenkins pipeline PASS — Bob's secure code allowed}
%     \label{fig:sprint4-jenkins-pass}
% \end{figure}

\textit{[Screenshot: Jenkins pipeline with green PASS stages]}

\subsection{Jenkins Console Output}

% TODO: INSERT SCREENSHOT OF JENKINS CONSOLE showing violation details
% \begin{figure}[H]
%     \centering
%     \includegraphics[width=0.9\textwidth]{figures/sprint4-console.png}
%     \caption{Sprint 4 — Jenkins console output showing violation details}
%     \label{fig:sprint4-console}
% \end{figure}

\textit{[Screenshot: Jenkins console with violation messages]}

\subsection{Dashboard After Demo}

% TODO: INSERT SCREENSHOT OF DASHBOARD after Alice + Bob scans
% \begin{figure}[H]
%     \centering
%     \includegraphics[width=0.9\textwidth]{figures/sprint4-dashboard-after.png}
%     \caption{Sprint 4 — Dashboard showing 2 scans, 50\% block rate after demo}
%     \label{fig:sprint4-dashboard-after}
% \end{figure}

\textit{[Screenshot: Dashboard showing 2 scans, 1 blocked, 1 allowed]}

% ============================================================
\section{Conclusion}
% ============================================================

Sprint 4 completed the DevSecOps Policy-as-Code microservice by integrating it into a real CI/CD pipeline. The Docker containerization ensures portable, reproducible deployment. The Jenkins pipeline integration proves that automated security scanning can be embedded seamlessly into existing CI/CD workflows. The Gitea setup demonstrates that the entire ecosystem can operate offline, making it suitable for demonstrations, education, and air-gapped environments.

The Alice and Bob scenario effectively illustrates the project's core value proposition: Alice's insecure infrastructure code is automatically detected and blocked before it can reach production, while Bob's compliant code passes through without delay. The dashboard provides real-time visibility into the team's security posture, and the detailed violation messages educate developers about security best practices.

With this sprint, the project achieved all of its initial objectives: 100 security policies, sub-50ms response time, multi-format support, educational feedback, a real-time dashboard, and complete CI/CD integration—all running entirely offline with zero cloud dependencies.
